\section{Problem Statement and Motivation}

The 2018 Danish drought caused agricultural losses exceeding DKK\nobreakspace 6\nobreakspace billion and exposed a structural gap in the national early warning landscape. Operational systems such as \textbf{HIP} (Hydrologisk Informations- og Prognosesystem) and the \textbf{DK model} deliver high resolution diagnostics but lack probabilistic, satellite integrated outlooks at sub-seasonal to seasonal lead times. Meanwhile, macroscale land surface models such as Variable Infiltration Capacity (\textbf{VIC}), the model targeted by the \textsc{DREAM} call, lack a prognostic groundwater state. In a country where more than 99\% of drinking water originates from aquifers, where water tables typically lie between 1 and 5\nobreakspace metres below surface, and where more than half of the agricultural area is sub-surface drained, this gap is not a detail but a first-order limitation. The recent finding by Forootan \textit{et\nobreakspace al.}\nobreakspace (2024) that the severity of the 2017 to 2021 European drought is systematically underestimated in non-assimilated \textbf{W3RA} confirms that closing this loop, through groundwater coupling, multi-source data assimilation, and hybrid physics and machine learning integration, is both timely and necessary. The \textsc{PhD} will deliver a satellite-fed VIC Denmark system producing a probabilistic four-pillar combined drought index and short- to seasonal-scale predictions of water dynamics.
